\section{Conclusion}
\label{sec:conclusion}

This paper evaluated 18 prompt-design variants across six open-weight,
instruction-tuned LLMs on the official EDOS sexism-detection benchmark
\cite{kirk2023semeval}, with paired significance testing, subtype-level
error analysis, order, few-shot-size, and demonstration-selection
sensitivity sweeps, and a head-to-head comparison against QLoRA, an
encoder, and a classical TF-IDF baseline. The clearest result in this
study cuts against prompting. Most of this study's prompt-design choices produce
large, statistically real effects (92 of 101 comparisons significant
under McNemar's test, 88 of 101 under the paired F1-diff bootstrap), and
phi3's largest single effect (+0.428 F1) is the biggest shift this study
puts through significance testing, though not the biggest shift in the
study overall: QLoRA fine-tuning moves the same model further still
(Sec.~\ref{sec:discussion:prompt-effects}). But
those effects operate entirely within a ceiling that a linear classifier
over word counts already clears for five of the six models, at a
fraction of the computational cost. Only qwen3 beats that classifier on
F1, and it does so at 4--16$\times$ the inference cost of every other
prompted model, while still losing to it on accuracy. Prompt design is a
real, measurable lever on this task, but the evidence here does not make
it a substitute for adapting a model, however simple, to the task's own
labeled data.

Two further findings deserve separate mention. Llama-3.1-8B's
interaction with chain-of-thought prompting is not the single-variant
curiosity it first appeared to be: every chain-of-thought variant this
model was evaluated on shows elevated generation failures, and the cause
splits cleanly depending on whether the prompt also asks for
aspect-by-aspect analysis, explicit safety refusal when it does not,
reasoning-token-budget exhaustion when it does
(Sec.~\ref{sec:results:llama-refusal}). Neither failure mode was visible
in the aggregate accuracy or F1 numbers alone; both surfaced only once
the raw generations were read directly. A model's headline score can
look ordinary while concealing a specific, systematic reason it is
wrong. Separately, prompt-\emph{component} order turned out to matter
far less than prior work on few-shot demonstration order would predict,
moving only one of six models' scores in either direction
(Sec.~\ref{sec:results:sensitivity}). In this study, which components a
prompt includes matters far more than the order they appear in.

For anyone building a sexism-detection or similar content-moderation
pipeline on an off-the-shelf instruction-tuned LLM, prompt engineering
is not wasted effort: it measurably changes what these models can do.
But it should be budgeted and evaluated as one option among several
rather than assumed to be the strongest one, and any pipeline that
treats a model's silence or refusal as an ordinary negative prediction
should monitor for that directly instead of trusting an aggregate score
to reveal it. The most
direct next steps are the ones this paper's limitations already name:
testing whether the same pattern holds for closed, commercial models,
and extending confidence-based scoring to chain-of-thought variants,
the one part of this study's own grid, including qwen3's best-performing
variant, that first-token AUC cannot reach at all.
